\section{Action and Reaction}

Although we have addressed the issue about the relationship between the spiritual authority and temporal power over and over, we may as well address it again with the hope and expectation that future commentators will offer more than links to other web sites. Here we are quoting Evola:

\begin{quotex}
In fact, such a thesis expresses the brahmanical-sacerdotal point of view of an Oriental, which, as we have already exposed and regardless of what Guenon may think, is but one of the possible points of view and by no means can pretend to have an absolute and exclusive value.
\end{quotex}


Obviously, Guenon is expressing the brahmanical point of view which, by its very nature is incomprehensible to other castes due to their differing inclinations and functions\footnote{\url{https://www.gornahoor.net/?p=1647}}. Of course, this ``point of view'' is more than an opinion, it is something that, according to Guenon and Tradition, can be
\emph{absolutely and exclusively} known.


\begin{quotex}
If that were the case, it would seem obvious that the Guelf thesis of the subordination of the State, as temporal power, to a spiritual authority monopolised by a sacerdotal caste, would be correct. But it is precisely the premise of this thesis that is not only erroneous but unacceptable because unilateral. Thus we must, in these circumstances, start again from zero.
\end{quotex}


Here is the source of Evola's confusion. Apparently he objects to the idea of ``subordination'', but that would be accurate only by confounding the notion of spiritual authority with temporal power. The brahmin caste does not exercise power in that way. The Traditional analogy is to regard them as the two wings of a bird, not as a relationship of subordination. That relationship is not at all unilateral, but accepted by both sides, hence it cannot be erroneous by Evola's argument. We have given the example, which Evola was certainly familiar with, of the Spartans and Romans at war\footnote{\url{https://www.gornahoor.net/?p=2024}}. The military caste waits for the priests to announce the auspicious moment. So even if we start again at zero, we still end up at the same place.

\begin{quotex}
This is for us a certain point which, among other things, clearly reveals the possibility of a spiritual significance of royalty, in virtue of which it can absorb and transcend in itself the sacerdotal function and be, beyond a temporal power, a spiritual authority.
\end{quotex}

Well, yes, royalty can combine both functions as we have documented in the Ancient City\footnote{\url{https://www.gornahoor.net/?p=2554}}. In the essay that Evola is critiquing, Guenon himself admits the possibility of the priest-king, since the two functions draw on a common higher principle. We take pains to point this out, as we recently did in regards to Buddha\footnote{\url{https://www.gornahoor.net/?p=3774}}. Closer to home, Christ has always been regarded as both Priest and King. So rather than refuting Guenon's ``point of view'', Evola talks like a man who just yesterday discovered sliced bread.

\begin{quotex}
From the purely metaphysical viewpoint, we do not speak of `principles' to know but of spiritual states to attain, of transcendent contacts to achieve owing to forces that, ultimately, do not cease to belong to the integral being of man.
\end{quotex}


Once again, we do not understand this objection. the attainment of spiritual states is identical to knowing those states; there is no distinction between knowing and being, a topic that we may have to address at another time. We disagree that all forces ultimately belong to the integral being of man. There are certainly forces that transcend the human state. We previously quoted Guenon:

\begin{quotex}
Action, no matter of what sort, cannot under any circumstances liberate from action; in other words it can only bear fruit within its own domain, which is that of human individuality. Thus it is not through action that it is possible to transcend individuality.
\end{quotex}


This is in perfect agreement with Evola, since Evola agrees that action is restricted to the integral being of man.

\begin{quotex}
In the metaphysical and concrete sense, `tradition' --- we have sufficiently repeated over and again --- is nothing more than the presence of such superior realisations as a continuity established from generation to generation by a chain of superior individualities. Wherever the royal tradition, defended by the warrior castes, has been correctly understood, it never had any other sense but that one. `Tradition' reduced to a doctrine, to a tradition of `teachings' and of `principles', is, in most cases, nothing but a \emph{caput mortuum} --- the `letter', much more than the `spirit'. And it is precisely on such grounds that `churches' flourish, whereas the experience of heroism and of self-control is something which is far more engaging for it to give place as often to such ambiguities and to such falls of tension.
\end{quotex}


This is a valid point, as far as it goes. When a tradition is reduced to a set of doctrines to be accepted in an extrinsic sense rather than an intrinsic knowing of them the tradition is dead. Unfortunately, the royal and warrior tradition dies with it. We see that in the decline of the Roman Empire, when the aristocracy no longer understood their own rites and traditions. Evola gives a strong defense of the Kshatriya perspective, bearing in mind that it is a perspective not a refutation of other perspectives. Since very few men, and fewer in the West, are called to seek Deliverance from all states, Evola's path is more appropriate as we have defended in \textit{The Cosmic Value of Human Action}\footnote{\url{https://www.gornahoor.net/?p=1856}}.

\begin{quotex}
Since Guenon understands correctly (insofar as it has always been the traditional teaching) that `knowledge', in the metaphysical sense, and `realisation' (action) are one and the same thing and inseparable elements in the simplicity of an act, Guenon should not have any difficulty admitting that action --- symbol of the warrior castes --- can constitute a path just as metaphysical as `knowledge'.
\end{quotex}


Once again, it seems that Evola is simply ignorant, since there is no argument about multiple paths. Guenon writes about the Path of the Gods (\emph{devayana}) and the Path of the ancestors (\emph{pitriyana}), where the goal of the former is the transcendence of all individual states. That is simply not the path of the warrior. In the Vedantic tradition, there is jnana yoga (knowledge), karma yoga (action) and bhakti yoga (devotion). So, obviously those are paths, but isn't it obvious to everyone that those who follow the path of action are not those who maintain knowledge of the tradition? In the Ancient City, the knowledge of the traditional rites was kept alive for centuries, but not necessarily by those who performed them.

\begin{quotex}
And perhaps should we remind Guenon of the Bhagavad Gita as an example of glorification of the warrior action that has nothing to do with the glorification of something material and temporal?
\end{quotex}


Well, Guenon refers to the Bhagavad Gita pretty often, himself. Nevertheless, did not Arjuna have to be convinced of the rightness of that warrior action? Krishna, the representative of the Brahmin caste, was Arjuna's teacher, so proper knowledge precedes right action. To regard the Gita as merely a text on the glorification of war rather than as a metaphysical teaching, is to greatly miss the point.

\begin{quotex}
And does not the very concept of `unmoved mover' used by Guenon as characterising the pure spiritual authority and defined as ``the thought of thought'', does it not refer to a modality of action, through this very notion of `mover'?
\end{quotex}


Of course, but what is the point? This \textbf{unmoved mover}, or Self, is precisely that which is unaffected by action, since it is unmoved. Yet as we pointed out in Salvation, Deliverance, Action\footnote{\url{https://www.gornahoor.net/?p=3780}}, even the enlightened being appears at a certain place and time, and hence appears to act. As we quoted Lama Govinda:

\begin{quotex}
He is no more a slave to law, but its master, because he has understood and realized it so profoundly. Through knowledge we master the law, and by mastering it, it ceases to be necessity, but becomes an instrument of real self-expression and spiritual freedom.
\end{quotex}


This is compatible with Evola's point of view as he has expressed it on many occasions. So the rupture he sees with Guenon is more apparent than real.

\begin{quotex}
It is not about a conflict between a spiritual authority and a rebellious temporal power but, on the contrary, a conflict between two distinct forms of authority equally spiritual and yet irreducible. Here again, we cannot repeat ourselves. Our readers know what we are talking about.
\end{quotex}


No, we don't know what he is talking about. Originally, the Ancient City was ruled by a priest-king. Then the temporal power revolted, stripping the kingly powers from the leader. This is historically documented in a work known to Evola (i.e., \emph{The Ancient City}).

\begin{quotex}
The cause cannot reside in the upper hand that, at a certain point, temporal power took spiritual authority. How can such a thing be possible in the first place? Should the hierarchy of which Guenon speaks thus be conceived as something so abstract, to the point of admitting that the superior does not also have the task of being the strongest? And if this were not the case, how could the inferior have imposed itself on the superior and thus paralysed the irresistible power of spiritual authority and supplanted it with temporal power?
\end{quotex}


Good question, and how could the warrior caste and landed aristocracy be taken over by merchants and bankers? And how could the plebes create a council in the heart of Rome, and later destroy the monarchies of France and Russia? Does it begin with the decline of spiritual authority, or the rebellion of lower castes who no longer support that authority? The idea that a warrior caste alone can unify an entire nation or Empire apart from a ``sacerdotal'' caste is impossible, both metaphysically and historically.

\paragraph{Conclusion}


Unfortunately, the overall tone of Evola's essay is polemical, rather than reasoned. Properly understood, there is ultimately little to choose between two points of view, which, as such, cannot refute each other, but should instead enhance each other. Evola is too eager to bring in irrelevancies such as the so-called ``sacerdotal'' caste, or lunar and solar traditions, a classification due to Bachofen, not any traditional writer. A case can be made for a Hermetic tradition and an heroic path. After all, in the Middle Ages both were living options and the warrior caste had its own initiatory path, separate from the ``sacerdotal'' caste. Nevertheless, the former accepted the spiritual authority (though not the ``power'') of the latter.

We can well understand the emotional power of Evola's vision. Every man can envision himself as a warrior, hero, or alchemist, but it is much more difficult to see himself as a yogi or priest. That is because he thinks he understands what it takes to be a warrior, but it is less clear what it means to be a yogi, or someone who has transcended all states. The temporal power can never assume spiritual authority, it can only lose it. It will then lose itself in the bargain.


\begin{center}* * *\end{center}

\begin{footnotesize}
\begin{sffamily}
\texttt{J on 2012-02-21 at 15:59 said:}

This is a useful response. I left the previous link with no argument because you supposedly eschew arguments. I figured people can read and differentiate/unify on their own.

My initial problem wasn't so much with your overall emphasis (else I wouldn't bother), but rather a minor point which probably comes down more to style and personality--particularly the incessant picking at others who apparently don't ``get it'', using what I think are demonstrably unfair tactics to dismiss a whole category of people who are taking an action-oriented approach involving something icky. I found it summarily dismissive and not really supported, since you failed to really identify of whom you were speaking and what makes their principles inherently incorrect. Lacking that, I had to presume it applies to all involved in anything labelled ``sex magic'', which renders a rather absurd reading, but since you grew it to include everyone who has sex at all, perhaps I still don't ``get it'', either. The whole attempt to bring in thelemites (without so labeling them, and apparently not at all dedicated to their definition of ``sex magic'') as some sort of counter-example was simply unnecessary and confusing.

Otherwise your points on metaphysics are well received and always bring about consideration, probably more than you would imagine, and I thank you for laying out a response to the issues raised in Evola's article.

Of course this talk of ``warriors'' may make the path of action seem more interesting, or more emotive than talk of priests. I wonder to what extent the warrior's experience of the nearness of death might transcend the priest's playing at it. I wonder so not because it entertains me to select a better, or to pretend at being a warrior, but because I don't think deliberately favoring one side or the other in these matters is really meaningful except to illustrate this or that point requiring the right perspective.

\hfill

\texttt{Cologero on 2012-02-21 at 23:12 said:}

The reason you don't ``get it'' is because what you are trying to get simply isn't there in the first place. If we posit a goal of reaching certain higher states, then of course we become interested in how to achieve them. If doing something icky can lead to higher states, or at least prepare oneself for them, then please explain how so. Obviously, it is not the ickyness in itself, but something superadded to it. Otherwise, the original criticism stands. So what is superadded?

I did indeed provide an example from Crowley which referred to his own description of the practice. Perhaps you know thelemites who differ from Crowley in that regard, but then the word has no meaning. As Crowley points out, the idea is to exhaust the body and mind, presumably to shut down the normal discursive thought processes in order to allow something else to happen in consciousness.

There is also a problem with self-deception. A few weeks ago, I read of an American TV show where the contenders for a prize had to drink donkey semen. My first reaction was to regard it as icky, but then I realized there are whole categories of people who may enjoy the taste of semen and regularly consume it. So it wouldn't really be a challenge for them.

For and understanding of action from different points of view, I can recommend these articles:

\textit{Qualifications for Initiation}\footnote{\url{http://www.gornahoor.net/?p=1161}}

\textit{Universality and the Heroic Life}\footnote{\url{http://www.gornahoor.net/?p=2885}}

\textit{The Three Trials}\footnote{\url{http://www.gornahoor.net/?p=1902}}

\textit{Builder of Worlds}\footnote{\url{http://www.gornahoor.net/?p=1884}}


\hfill

\texttt{logres on 2012-02-22 at 00:44 said:}

Where does the ``prophet'' fit in?

\hfill

\texttt{Losang Shenphen on 2012-02-23 at 12:35 said:}

Excellent observation, Mr Hoo, Gornahoor is quite amusing, one big joke in fact. It should only be read through a monocle.

\hfill

\texttt{HOO on 2012-02-23 at 13:29 said:}

Losang Shenphen, why do you say this? And what is the use of a monocle when one has perfect vision?

\hfill

\texttt{HOO on 2012-02-23 at 18:17 said:}

I'm not sure if anybody meant so, but to be clear, Traditional Sexual-Magic is not merely symbolic. Traditional sexual-magic involves the manipulation of energy (``subtle and vital currents (vayus)'') latent within the body. Nor is anything symbolic, when it is regarding a true symbol, merely so; as symbols are ``\,''magic, menial, psychic and moral operations'' awakening new notions, ideas, sentiments and aspirations.'' And they are always reflected in the organism.\\ Regarding inefficient sexual-rituals, if those include those of Crowley or not; why would we worry that others fool themselves, enjoy themselves or suffer? Pashu does as pashu is. Cologero ''To regard the Gita as MERELY a text on the glorification of war rather than as a metaphysical teaching''\\ Evola nowhere does this.

\hfill

\texttt{apeiron on 2012-02-23 at 21:20 said:}

Evola has essentially understood that the two were once one that converged with each other (both temporal and spiritual power, the two-headed alchemical androgyne, Priest-King, Sacral Imperator).

This makes the most sense, as demonstrated linearly within an historical manifestation; all emanations beyond the first point are fractals that project a material reflection of an absolute paradisaical primordiality. Call them lunar shadows if you will, but Evola ultimately recognized their unified existence of the unmanifested absolute being. The trouble begins when someone argues for the legitimacy of one power above the other, when Evola simply makes the argument that they are both as legitimate and converged at the origins into a healthy symbiosis. Will, intellect and action united, this is the meaning of the Cross.

Once again, these breakaways(castes, ways of being, etc) if historically demonstrated, will fracture further away from the center and eventually revolt against each other into lower forms. Rule no.1 one: stay centered, avoid fall. Which is fundamentally why, origins, properly understood, are so integral to Tradition.

\hfill

\texttt{Cologero on 2012-02-23 at 22:34 said:}

Yes, Hoo, I added the word MERELY. In all cases, the open question is exactly what is beyond the material and temporal?

So to return to sex magic, what is beyond the material and temporal? You imply subtle energy is within the body. Let us look at what Guenon writes:

The gross state is corporeal existence.\\ The subtle state includes extra-corporeal modalities of the human being as well as all other individual states.\\ These two terms are not truly symmetrical and cannot even have any common measure.

Hence, there are no subtle energies latent within the body, which would be a contradiction. So the subtle energies are to be found elsewhere than in the body. So it is far from clear how a physical act can awaken the awareness of such subtle energies, at least in any mechanical way. Of course, physical trials may provide the opportunity to confront one's fears, weaknesses, lusts and so on, and to develop the opposite qualities of courage, strength, fidelity. But this is far from automatic.

We have yet to hear from anyone regarding the specifics of sex magic in this regard. Also, oddly enough, no one has bothered to respond to this question: if we accept alchemical descriptions as symbolic, then why not the same for sex magic?

We have received almost no comments on Serrano's
\textit{The Test}\footnote{\url{http://www.gornahoor.net/?p=70}}. Here is some homework:

practice sex magic with total detachment and no expectation; then report back on the results.

As for the fools, we have great compassion for them and endeavour to save them years of fruitless searching. Also, fools manage to teach us a lot. Isn't that the real meaning of the Tantras, the practice for the Kali Yuga? Specifically, that enlightenment is no longer constrained to certain paths and practices, but may arise from formerly forbidden sources, such as fools, women, or outcastes.

\hfill

\texttt{Cologero on 2012-02-23 at 23:19 said:}

Yes, with a caveat. Guenon and Evola write of spiritual
\emph{authority}, not power. Resentment is the sign of the kali yuga as well as the idea that one caste is \emph{over} another. Rather, it is a question of their respective functions or spheres of influence.

\hfill


\end{sffamily}
\end{footnotesize}

